% =====================================================================
%  Ejercicio20CorreccionMaxHeapify.tex
%  ========================
%
%
%  Descripción:
%  ------------
%  Ejercicio 20 de la hoja de cotas asintóticas, recurrencias y
%  montículos, resuelto paso a paso. Un ejercicio por documento.
%
%
%  Considerations:
%  ------------
%  - Compila con pdflatex sin shell-escape. Dos pasadas, o latexmk
%  - Solo declara las convenciones que este ejercicio usa
%  - Sin título grande y sin pie de página: el encabezado de un
%    renglón deja la hoja entera para el desarrollo
%
%
%  Metadata:
%  ----------
%   * Author : zxxz6 (Bryan Violante Arriaga)
%   * Version: 1.0.0
%
%
%  History:
%  ------------
%  Author      Date            Description
%  zxxz6       24/09/2026      Creation
%
%
% =====================================================================
\documentclass[11pt,letterpaper]{article}

\newcommand{\materia}{Hoja de recurrencias y montículos}
\newcommand{\profesor}{Dr. José Martínez Carranza}
\newcommand{\periodo}{Otoño 2026}
\newcommand{\alumno}{Bryan Ezequiel Violante Arriaga}

% =====================================================================
%  Preambulo.tex
%  ========================
%
%
%  Descripción:
%  ------------
%  Preámbulo compartido de los apuntes de clase: paquetes, estilo de
%  página, entornos de teorema, las cuatro cajas de color, el estilo de
%  listings y el pseudocódigo en español. Se carga con \input desde
%  cada documento en lugar de copiarlo, para que un arreglo aquí llegue
%  a todas las libretas de una vez.
%
%
%  Considerations:
%  ------------
%  - Se carga con \input y ruta relativa, no con \usepackage, porque
%    \usepackage solo busca en la carpeta del documento y en el árbol
%    de TeX. Un .sty instalado en TEXMFHOME evitaría la ruta, pero vive
%    fuera del repo y una clona nueva no compilaría
%  - Los cuatro datos de la materia se declaran en el documento ANTES
%    del \input. hyperref expande \materia al cargarse, así que si no
%    existe todavía la compilación truena
%  - Cuando la materia la imparte más de un profesor, los nombres van
%    en un solo \profesor separados por \\. La portadilla los imprime
%    dentro de un center, que es donde \\ sí corta renglón. Repetir
%    \newcommand{\profesor} no es opción: truena con "already defined"
%  - algorithmicx no tiene opción de idioma y babel no lo traduce, así
%    que las palabras clave del pseudocódigo se renombran a mano
%  - Los nombres de los entornos van sin acento (\begin{definicion})
%    aunque el rótulo impreso sí lo lleve
%  - babel se carga con es-noshorthands porque en español el " queda
%    activo y se come el texto que le sigue, además de romper listings
%
%
%  Metadata:
%  ----------
%   * Author : zxxz6 (Bryan Violante Arriaga)
%   * Version: 1.7.0
%
%
%  History:
%  ------------
%  Author      Date            Description
%  zxxz6       24/09/2026      El nombre del alumno en el encabezado
%  zxxz6       24/09/2026      Encabezado de una hoja de un ejercicio
%  zxxz6       23/09/2026      Macros de los documentos de soluciones
%  zxxz6       30/08/2026      Macros de lim, limsup y liminf sobre n
%  zxxz6       25/08/2026      Macros de Omega, Theta, o y omega chicas
%  zxxz6       20/08/2026      Cargue pgfplots para graficas de datos
%  zxxz6       20/08/2026      Cargue tikz para diagramas de conjuntos
%  zxxz6       19/08/2026      Documente \profesor con varios nombres
%  zxxz6       19/08/2026      Agregue la caja Idea Random
%  zxxz6       18/08/2026      Creation
%
%
% =====================================================================
%  USO
%  ---
%  \documentclass[11pt,letterpaper]{article}
%  \newcommand{\materia}{...}      % los cuatro datos van arriba
%  \newcommand{\profesor}{...}
%  \newcommand{\periodo}{...}
%  \newcommand{\alumno}{...}
%  % =====================================================================
%  Preambulo.tex
%  ========================
%
%
%  Descripción:
%  ------------
%  Preámbulo compartido de los apuntes de clase: paquetes, estilo de
%  página, entornos de teorema, las cuatro cajas de color, el estilo de
%  listings y el pseudocódigo en español. Se carga con \input desde
%  cada documento en lugar de copiarlo, para que un arreglo aquí llegue
%  a todas las libretas de una vez.
%
%
%  Considerations:
%  ------------
%  - Se carga con \input y ruta relativa, no con \usepackage, porque
%    \usepackage solo busca en la carpeta del documento y en el árbol
%    de TeX. Un .sty instalado en TEXMFHOME evitaría la ruta, pero vive
%    fuera del repo y una clona nueva no compilaría
%  - Los cuatro datos de la materia se declaran en el documento ANTES
%    del \input. hyperref expande \materia al cargarse, así que si no
%    existe todavía la compilación truena
%  - Cuando la materia la imparte más de un profesor, los nombres van
%    en un solo \profesor separados por \\. La portadilla los imprime
%    dentro de un center, que es donde \\ sí corta renglón. Repetir
%    \newcommand{\profesor} no es opción: truena con "already defined"
%  - algorithmicx no tiene opción de idioma y babel no lo traduce, así
%    que las palabras clave del pseudocódigo se renombran a mano
%  - Los nombres de los entornos van sin acento (\begin{definicion})
%    aunque el rótulo impreso sí lo lleve
%  - babel se carga con es-noshorthands porque en español el " queda
%    activo y se come el texto que le sigue, además de romper listings
%
%
%  Metadata:
%  ----------
%   * Author : zxxz6 (Bryan Violante Arriaga)
%   * Version: 1.7.0
%
%
%  History:
%  ------------
%  Author      Date            Description
%  zxxz6       24/09/2026      El nombre del alumno en el encabezado
%  zxxz6       24/09/2026      Encabezado de una hoja de un ejercicio
%  zxxz6       23/09/2026      Macros de los documentos de soluciones
%  zxxz6       30/08/2026      Macros de lim, limsup y liminf sobre n
%  zxxz6       25/08/2026      Macros de Omega, Theta, o y omega chicas
%  zxxz6       20/08/2026      Cargue pgfplots para graficas de datos
%  zxxz6       20/08/2026      Cargue tikz para diagramas de conjuntos
%  zxxz6       19/08/2026      Documente \profesor con varios nombres
%  zxxz6       19/08/2026      Agregue la caja Idea Random
%  zxxz6       18/08/2026      Creation
%
%
% =====================================================================
%  USO
%  ---
%  \documentclass[11pt,letterpaper]{article}
%  \newcommand{\materia}{...}      % los cuatro datos van arriba
%  \newcommand{\profesor}{...}
%  \newcommand{\periodo}{...}
%  \newcommand{\alumno}{...}
%  % =====================================================================
%  Preambulo.tex
%  ========================
%
%
%  Descripción:
%  ------------
%  Preámbulo compartido de los apuntes de clase: paquetes, estilo de
%  página, entornos de teorema, las cuatro cajas de color, el estilo de
%  listings y el pseudocódigo en español. Se carga con \input desde
%  cada documento en lugar de copiarlo, para que un arreglo aquí llegue
%  a todas las libretas de una vez.
%
%
%  Considerations:
%  ------------
%  - Se carga con \input y ruta relativa, no con \usepackage, porque
%    \usepackage solo busca en la carpeta del documento y en el árbol
%    de TeX. Un .sty instalado en TEXMFHOME evitaría la ruta, pero vive
%    fuera del repo y una clona nueva no compilaría
%  - Los cuatro datos de la materia se declaran en el documento ANTES
%    del \input. hyperref expande \materia al cargarse, así que si no
%    existe todavía la compilación truena
%  - Cuando la materia la imparte más de un profesor, los nombres van
%    en un solo \profesor separados por \\. La portadilla los imprime
%    dentro de un center, que es donde \\ sí corta renglón. Repetir
%    \newcommand{\profesor} no es opción: truena con "already defined"
%  - algorithmicx no tiene opción de idioma y babel no lo traduce, así
%    que las palabras clave del pseudocódigo se renombran a mano
%  - Los nombres de los entornos van sin acento (\begin{definicion})
%    aunque el rótulo impreso sí lo lleve
%  - babel se carga con es-noshorthands porque en español el " queda
%    activo y se come el texto que le sigue, además de romper listings
%
%
%  Metadata:
%  ----------
%   * Author : zxxz6 (Bryan Violante Arriaga)
%   * Version: 1.7.0
%
%
%  History:
%  ------------
%  Author      Date            Description
%  zxxz6       24/09/2026      El nombre del alumno en el encabezado
%  zxxz6       24/09/2026      Encabezado de una hoja de un ejercicio
%  zxxz6       23/09/2026      Macros de los documentos de soluciones
%  zxxz6       30/08/2026      Macros de lim, limsup y liminf sobre n
%  zxxz6       25/08/2026      Macros de Omega, Theta, o y omega chicas
%  zxxz6       20/08/2026      Cargue pgfplots para graficas de datos
%  zxxz6       20/08/2026      Cargue tikz para diagramas de conjuntos
%  zxxz6       19/08/2026      Documente \profesor con varios nombres
%  zxxz6       19/08/2026      Agregue la caja Idea Random
%  zxxz6       18/08/2026      Creation
%
%
% =====================================================================
%  USO
%  ---
%  \documentclass[11pt,letterpaper]{article}
%  \newcommand{\materia}{...}      % los cuatro datos van arriba
%  \newcommand{\profesor}{...}
%  \newcommand{\periodo}{...}
%  \newcommand{\alumno}{...}
%  \input{../../templates/Preambulo.tex}
%  \begin{document}
%  \portadilla
%
%  Con dos o mas profesores se declara UN solo \profesor y los nombres
%  se separan con \\, uno por renglon:
%
%  \newcommand{\profesor}{Dra. Fulana de Tal \\
%                         Dr. Mengano de Cual}
% =====================================================================


% ==================== IDIOMA Y CODIFICACIÓN ====================
\usepackage[utf8]{inputenc}
\usepackage[T1]{fontenc}
\usepackage[spanish,es-tabla,es-noshorthands]{babel}
\usepackage{lmodern}
\usepackage{microtype}

% ==================== PÁGINA ====================
\usepackage[letterpaper,margin=2.3cm,headheight=14pt]{geometry}
\usepackage{parskip}
\usepackage{fancyhdr}
\usepackage{enumitem}

% ==================== MATEMÁTICAS ====================
\usepackage{amsmath,amssymb,amsthm,mathtools}

% ==================== FIGURAS, TABLAS, CÓDIGO ====================
\usepackage{graphicx}
\usepackage{booktabs}
\usepackage{array}
% float da el especificador [H], que clava una tabla donde se
% escribio. Una traza que LaTeX manda dos paginas adelante deja
% de leerse junto al paso que la produjo
\usepackage{float}
\usepackage{xcolor}
% tikz ya viene arrastrado por tcolorbox[most], pero se carga
% explicito para no depender de ese detalle; la libreria babel
% protege a tikz de los caracteres activos del español
\usepackage{tikz}
\usetikzlibrary{babel}
% pgfplots dibuja graficas de datos (ejes, escalas log) sobre tikz
\usepackage{pgfplots}
\pgfplotsset{compat=1.18}
\usepackage{listings}
\usepackage[ruled,section]{algorithm}
\usepackage{algpseudocode}
\usepackage[most]{tcolorbox}
% soul da \hl, el unico resaltado que parte renglon; \colorbox no
\usepackage{soul}
\usepackage{hyperref}

% =====================================================================
%  ESTILO
% =====================================================================

% --- encabezado y pie de cada página ---
\pagestyle{fancy}
\fancyhf{}
\fancyhead[L]{\small\textsc{\materia}}
\fancyhead[R]{\small\periodo}
\fancyfoot[C]{\small\thepage}
\renewcommand{\headrulewidth}{0.4pt}

% --- enlaces ---
\hypersetup{
  colorlinks=true,
  linkcolor=black,
  urlcolor=blue,
  citecolor=black,
  pdftitle={Apuntes de \materia},
  pdfauthor={\alumno}
}

% --- paleta ---
\definecolor{acento}{HTML}{1F4E79}
\definecolor{fondoIdea}{HTML}{EAF2FB}
\definecolor{fondoOjo}{HTML}{FDF0E6}
\definecolor{fondoPend}{HTML}{F2F0F7}
\definecolor{comentario}{HTML}{5A7D5A}
\definecolor{cadena}{HTML}{A03030}
\definecolor{fondoCodigo}{HTML}{F7F7F7}
\definecolor{fondoResultado}{HTML}{EAF5EC}
\definecolor{marcador}{HTML}{FFF3A3}

% --- entornos de teorema; numeran por sección, o sea por clase ---
\theoremstyle{definition}
\newtheorem{definicion}{Definición}[section]
\newtheorem{ejemplo}[definicion]{Ejemplo}
\newtheorem{ejercicio}[definicion]{Ejercicio}

\theoremstyle{plain}
\newtheorem{teorema}[definicion]{Teorema}
\newtheorem{lema}[definicion]{Lema}
\newtheorem{corolario}[definicion]{Corolario}
\newtheorem{proposicion}[definicion]{Proposición}

\theoremstyle{remark}
\newtheorem*{observacion}{Observación}
\newtheorem*{quick_sum}{Resumen rápido}

% --- cajas de color para lo que se relee antes del examen ---
\newtcolorbox{idea}[1][]{
  colback=fondoIdea, colframe=acento, boxrule=0.6pt,
  arc=2pt, left=6pt, right=6pt, top=5pt, bottom=5pt,
  title={Idea clave}, fonttitle=\bfseries\small, #1
}
\newtcolorbox{ojo}[1][]{
  colback=fondoOjo, colframe=orange!70!black, boxrule=0.6pt,
  arc=2pt, left=6pt, right=6pt, top=5pt, bottom=5pt,
  title={Ojito...}, fonttitle=\bfseries\small, #1
}
\newtcolorbox{pendiente}[1][]{
  colback=fondoPend, colframe=violet!60!black, boxrule=0.6pt,
  arc=2pt, left=6pt, right=6pt, top=5pt, bottom=5pt,
  title={Pendiente por resolver}, fonttitle=\bfseries\small, #1
}
\newtcolorbox{idear}[1][]{
  colback=fondoPend, colframe=pink!60!black, boxrule=0.6pt,
  arc=2pt, left=6pt, right=6pt, top=5pt, bottom=5pt,
  title={Idea Random}, fonttitle=\bfseries\small, #1
}

% --- el resultado al que llego un ejercicio, para que se vea de
%     lejos al hojear la hoja antes del examen ---
\newtcolorbox{resultado}[1][]{
  colback=fondoResultado, colframe=green!45!black, boxrule=0.6pt,
  arc=2pt, left=6pt, right=6pt, top=5pt, bottom=5pt,
  title={Resultado}, fonttitle=\bfseries\small, #1
}

% --- código ---
\lstdefinestyle{codigo}{
  backgroundcolor=\color{fondoCodigo},
  basicstyle=\ttfamily\footnotesize,
  keywordstyle=\color{acento}\bfseries,
  commentstyle=\color{comentario}\itshape,
  stringstyle=\color{cadena},
  numbers=left, numberstyle=\tiny\color{gray}, numbersep=8pt,
  frame=single, rulecolor=\color{gray!40},
  breaklines=true, showstringspaces=false,
  tabsize=4, captionpos=b,
  extendedchars=true, inputencoding=utf8,
  literate={á}{{\'a}}1 {é}{{\'e}}1 {í}{{\'i}}1
           {ó}{{\'o}}1 {ú}{{\'u}}1 {ñ}{{\~n}}1
           {Ü}{{\"U}}1 {ü}{{\"u}}1 {¿}{{?`}}1
           {¡}{{!`}}1
}
\lstset{style=codigo}
\renewcommand{\lstlistingname}{Código}
\renewcommand{\lstlistlistingname}{Índice de código}
% --- pseudocódigo; algorithmicx arma el cuerpo y algorithm el
%     flotante que lo numera y lo pone en el índice. Las palabras
%     clave vienen en inglés de fábrica y babel no las toca, así que
%     se renombran una por una ---
\floatname{algorithm}{Algoritmo}
\renewcommand{\listalgorithmname}{Índice de algoritmos}

\algrenewcommand\algorithmicrequire{\textbf{Entrada:}}
\algrenewcommand\algorithmicensure{\textbf{Salida:}}
\algrenewcommand\algorithmicend{\textbf{fin}}
\algrenewcommand\algorithmicif{\textbf{si}}
\algrenewcommand\algorithmicthen{\textbf{entonces}}
\algrenewcommand\algorithmicelse{\textbf{si no}}
\algrenewcommand\algorithmicfor{\textbf{para}}
\algrenewcommand\algorithmicforall{\textbf{para todo}}
\algrenewcommand\algorithmicdo{\textbf{hacer}}
\algrenewcommand\algorithmicwhile{\textbf{mientras}}
\algrenewcommand\algorithmicrepeat{\textbf{repetir}}
\algrenewcommand\algorithmicuntil{\textbf{hasta que}}
\algrenewcommand\algorithmicloop{\textbf{ciclo}}
\algrenewcommand\algorithmicprocedure{\textbf{procedimiento}}
\algrenewcommand\algorithmicfunction{\textbf{función}}
\algrenewcommand\algorithmicreturn{\textbf{devolver}}
\algrenewcomment[1]{\hfill{\color{comentario}\itshape // #1}}

% Las dos que algorithmicx no trae. El \ del final es obligatorio: sin
% el TeX se come el espacio y sale "hastan-1" en vez de "hasta n-1"
\algnewcommand\Hasta{\textbf{hasta}\ }
\algnewcommand\Intercambiar{\textbf{intercambiar}\ }

% --- abreviaturas que se usan a cada rato ---
\newcommand{\N}{\mathbb{N}}
\newcommand{\Z}{\mathbb{Z}}
\newcommand{\Q}{\mathbb{Q}}
\newcommand{\R}{\mathbb{R}}
\newcommand{\C}{\mathbb{C}}
\newcommand{\abs}[1]{\left\lvert #1 \right\rvert}
\newcommand{\norm}[1]{\left\lVert #1 \right\rVert}
\newcommand{\set}[1]{\left\{ #1 \right\}}
\newcommand{\Ord}[1]{\mathcal{O}\!\left(#1\right)}
% Mayuscula la notacion grande, minuscula la chica
\newcommand{\Omg}[1]{\Omega\!\left(#1\right)}
\newcommand{\Tht}[1]{\Theta\!\left(#1\right)}
\newcommand{\ord}[1]{o\!\left(#1\right)}
\newcommand{\omg}[1]{\omega\!\left(#1\right)}
% Los tres limites sobre n, que salen en cada renglon del criterio
% del cociente para la notacion asintotica
\newcommand{\LM}{\lim_{n \to \infty}}
\newcommand{\LS}{\limsup_{n \to \infty}}
\newcommand{\LI}{\liminf_{n \to \infty}}

% --- encabezado de cada sesión: título a la izquierda, fecha a la
%     derecha. El argumento corto es el que sale en el índice ---
\newcommand{\clase}[2]{%
  \section[#1]{#1 \hfill \normalfont\small\textit{#2}}%
}

% --- portadilla y tabla de contenidos. Se llama una vez, justo
%     despues de \begin{document}. El renglon del profesor va dentro
%     del center, asi que \profesor puede traer varios nombres
%     separados por \\ y salen apilados y centrados sin tocar nada ---
\newcommand{\portadilla}{%
  \begin{center}
    {\LARGE\bfseries\materia\par}
    \vspace{0.4em}
    {\large \profesor \par}
    \vspace{0.2em}
    {\small \periodo\quad-\quad\alumno\par}
    \vspace{0.3em}
    \rule{\linewidth}{0.5pt}
  \end{center}
  \vspace{0.5em}
  \tableofcontents
  \vspace{1em}
  \rule{\linewidth}{0.5pt}%
}

% =====================================================================
%  DOCUMENTOS DE SOLUCIONES
%  Lo que usa una tanda de ejercicios resueltos. Una libreta no toca
%  nada de aqui, pero vive en el preambulo compartido para no tener
%  dos archivos que arreglar cuando cambie el estilo.
% =====================================================================

% --- encabezado compacto: titulo, subtitulo y regla. Sustituye a
%     \portadilla cuando el documento es una tanda de ejercicios: sin
%     indice, porque se hojea de corrido, no se navega ---
\newcommand{\encabezadoSoluciones}[2]{%
  % El paquete algorithm se carga con la opcion [section] para que en
  % una libreta el pseudocodigo se numere por clase. Aqui no hay
  % secciones, asi que esa numeracion sale como "Algoritmo 0.1". Se
  % corrige al vuelo, y solo en los documentos que llaman a este
  % encabezado: la libreta no se entera.
  \renewcommand{\thealgorithm}{\arabic{algorithm}}%
  \begin{center}
    {\LARGE\bfseries #1\par}
    \vspace{0.3em}
    {\small\itshape #2\par}
    \vspace{0.5em}
    \rule{\linewidth}{0.5pt}
  \end{center}
  \vspace{0.3em}%
}

% --- hoja de un solo ejercicio: un renglon chico arriba con el
%     titulo y su numero, y ya. Sin titulo grande y sin pie de pagina,
%     porque en una hoja de dos paginas cada centimetro que se va en
%     adornos es un paso del desarrollo que no cabe ---
\newcommand{\encabezadoEjercicio}[3]{%
  % Misma correccion que en \encabezadoSoluciones: sin secciones, el
  % pseudocodigo saldria numerado "0.1"
  \renewcommand{\thealgorithm}{\arabic{algorithm}}%
  \fancyhf{}%
  \fancyhead[L]{\small\textsc{#1} (ejercicio #2 de #3)}%
  % El nombre va a la derecha porque las hojas de respuestas se
  % entregan y el profesor pide que cada una lo lleve
  \fancyhead[R]{\small\alumno}%
  \renewcommand{\headrulewidth}{0.4pt}%
  \pagestyle{fancy}%
  \thispagestyle{fancy}%
}

% --- bloque tematico, numerado con romanos: I, II, III. Agrupa los
%     ejercicios que comparten tema. No dibuja regla: la del
%     \problema que sigue ya separa el titulo del primer ejercicio,
%     y dos rayas seguidas dejan un hueco feo ---
\newcounter{parte}
\renewcommand{\theparte}{\Roman{parte}}
\newcommand{\parte}[1]{%
  \bigskip
  \refstepcounter{parte}%
  \noindent{\large\bfseries \theparte. #1\par}%
}

% --- un ejercicio resuelto. Se numera solo y corrido a lo largo del
%     documento, sin reiniciar en cada parte, para poder decir "el 14"
%     sin aclarar de que bloque ---
\newcounter{problema}
\newcommand{\problema}[1]{%
  \bigskip
  \noindent\rule{\linewidth}{0.4pt}\par
  \vspace{0.3em}
  \refstepcounter{problema}%
  \noindent{\bfseries Ejercicio \theproblema\ --- #1\par}
  \vspace{0.3em}%
}

% --- rotulo del paso que sigue: Caso base, Hipotesis inductiva, Paso
%     inductivo. El texto sigue en el mismo renglon ---
\newcommand{\paso}[1]{\par\vspace{0.3em}\noindent\textbf{#1}\ }

% --- la respuesta de un ejercicio, centrada y en un cuadro. El
%     argumento va en modo matematico, sin los delimitadores ---
\newcommand{\respuesta}[1]{%
  \[
    \boxed{\;#1\;}
  \]
}

% --- resaltado inline, como pasarle el plumon a media frase. No
%     admite matematicas adentro: soul rompe con $...$, asi que una
%     formula que hay que destacar va en \respuesta ---
\sethlcolor{marcador}
\newcommand{\resaltar}[1]{\hl{#1}}

% --- justificacion al final de un renglon de \align, en chiquito y
%     entre parentesis, para decir de donde salio el paso ---
\newcommand{\razon}[1]{\quad\text{\small\itshape (#1)}}

%############################ END OF PREAMBULO.TEX #############################
%###############################################################################

%  \begin{document}
%  \portadilla
%
%  Con dos o mas profesores se declara UN solo \profesor y los nombres
%  se separan con \\, uno por renglon:
%
%  \newcommand{\profesor}{Dra. Fulana de Tal \\
%                         Dr. Mengano de Cual}
% =====================================================================


% ==================== IDIOMA Y CODIFICACIÓN ====================
\usepackage[utf8]{inputenc}
\usepackage[T1]{fontenc}
\usepackage[spanish,es-tabla,es-noshorthands]{babel}
\usepackage{lmodern}
\usepackage{microtype}

% ==================== PÁGINA ====================
\usepackage[letterpaper,margin=2.3cm,headheight=14pt]{geometry}
\usepackage{parskip}
\usepackage{fancyhdr}
\usepackage{enumitem}

% ==================== MATEMÁTICAS ====================
\usepackage{amsmath,amssymb,amsthm,mathtools}

% ==================== FIGURAS, TABLAS, CÓDIGO ====================
\usepackage{graphicx}
\usepackage{booktabs}
\usepackage{array}
% float da el especificador [H], que clava una tabla donde se
% escribio. Una traza que LaTeX manda dos paginas adelante deja
% de leerse junto al paso que la produjo
\usepackage{float}
\usepackage{xcolor}
% tikz ya viene arrastrado por tcolorbox[most], pero se carga
% explicito para no depender de ese detalle; la libreria babel
% protege a tikz de los caracteres activos del español
\usepackage{tikz}
\usetikzlibrary{babel}
% pgfplots dibuja graficas de datos (ejes, escalas log) sobre tikz
\usepackage{pgfplots}
\pgfplotsset{compat=1.18}
\usepackage{listings}
\usepackage[ruled,section]{algorithm}
\usepackage{algpseudocode}
\usepackage[most]{tcolorbox}
% soul da \hl, el unico resaltado que parte renglon; \colorbox no
\usepackage{soul}
\usepackage{hyperref}

% =====================================================================
%  ESTILO
% =====================================================================

% --- encabezado y pie de cada página ---
\pagestyle{fancy}
\fancyhf{}
\fancyhead[L]{\small\textsc{\materia}}
\fancyhead[R]{\small\periodo}
\fancyfoot[C]{\small\thepage}
\renewcommand{\headrulewidth}{0.4pt}

% --- enlaces ---
\hypersetup{
  colorlinks=true,
  linkcolor=black,
  urlcolor=blue,
  citecolor=black,
  pdftitle={Apuntes de \materia},
  pdfauthor={\alumno}
}

% --- paleta ---
\definecolor{acento}{HTML}{1F4E79}
\definecolor{fondoIdea}{HTML}{EAF2FB}
\definecolor{fondoOjo}{HTML}{FDF0E6}
\definecolor{fondoPend}{HTML}{F2F0F7}
\definecolor{comentario}{HTML}{5A7D5A}
\definecolor{cadena}{HTML}{A03030}
\definecolor{fondoCodigo}{HTML}{F7F7F7}
\definecolor{fondoResultado}{HTML}{EAF5EC}
\definecolor{marcador}{HTML}{FFF3A3}

% --- entornos de teorema; numeran por sección, o sea por clase ---
\theoremstyle{definition}
\newtheorem{definicion}{Definición}[section]
\newtheorem{ejemplo}[definicion]{Ejemplo}
\newtheorem{ejercicio}[definicion]{Ejercicio}

\theoremstyle{plain}
\newtheorem{teorema}[definicion]{Teorema}
\newtheorem{lema}[definicion]{Lema}
\newtheorem{corolario}[definicion]{Corolario}
\newtheorem{proposicion}[definicion]{Proposición}

\theoremstyle{remark}
\newtheorem*{observacion}{Observación}
\newtheorem*{quick_sum}{Resumen rápido}

% --- cajas de color para lo que se relee antes del examen ---
\newtcolorbox{idea}[1][]{
  colback=fondoIdea, colframe=acento, boxrule=0.6pt,
  arc=2pt, left=6pt, right=6pt, top=5pt, bottom=5pt,
  title={Idea clave}, fonttitle=\bfseries\small, #1
}
\newtcolorbox{ojo}[1][]{
  colback=fondoOjo, colframe=orange!70!black, boxrule=0.6pt,
  arc=2pt, left=6pt, right=6pt, top=5pt, bottom=5pt,
  title={Ojito...}, fonttitle=\bfseries\small, #1
}
\newtcolorbox{pendiente}[1][]{
  colback=fondoPend, colframe=violet!60!black, boxrule=0.6pt,
  arc=2pt, left=6pt, right=6pt, top=5pt, bottom=5pt,
  title={Pendiente por resolver}, fonttitle=\bfseries\small, #1
}
\newtcolorbox{idear}[1][]{
  colback=fondoPend, colframe=pink!60!black, boxrule=0.6pt,
  arc=2pt, left=6pt, right=6pt, top=5pt, bottom=5pt,
  title={Idea Random}, fonttitle=\bfseries\small, #1
}

% --- el resultado al que llego un ejercicio, para que se vea de
%     lejos al hojear la hoja antes del examen ---
\newtcolorbox{resultado}[1][]{
  colback=fondoResultado, colframe=green!45!black, boxrule=0.6pt,
  arc=2pt, left=6pt, right=6pt, top=5pt, bottom=5pt,
  title={Resultado}, fonttitle=\bfseries\small, #1
}

% --- código ---
\lstdefinestyle{codigo}{
  backgroundcolor=\color{fondoCodigo},
  basicstyle=\ttfamily\footnotesize,
  keywordstyle=\color{acento}\bfseries,
  commentstyle=\color{comentario}\itshape,
  stringstyle=\color{cadena},
  numbers=left, numberstyle=\tiny\color{gray}, numbersep=8pt,
  frame=single, rulecolor=\color{gray!40},
  breaklines=true, showstringspaces=false,
  tabsize=4, captionpos=b,
  extendedchars=true, inputencoding=utf8,
  literate={á}{{\'a}}1 {é}{{\'e}}1 {í}{{\'i}}1
           {ó}{{\'o}}1 {ú}{{\'u}}1 {ñ}{{\~n}}1
           {Ü}{{\"U}}1 {ü}{{\"u}}1 {¿}{{?`}}1
           {¡}{{!`}}1
}
\lstset{style=codigo}
\renewcommand{\lstlistingname}{Código}
\renewcommand{\lstlistlistingname}{Índice de código}
% --- pseudocódigo; algorithmicx arma el cuerpo y algorithm el
%     flotante que lo numera y lo pone en el índice. Las palabras
%     clave vienen en inglés de fábrica y babel no las toca, así que
%     se renombran una por una ---
\floatname{algorithm}{Algoritmo}
\renewcommand{\listalgorithmname}{Índice de algoritmos}

\algrenewcommand\algorithmicrequire{\textbf{Entrada:}}
\algrenewcommand\algorithmicensure{\textbf{Salida:}}
\algrenewcommand\algorithmicend{\textbf{fin}}
\algrenewcommand\algorithmicif{\textbf{si}}
\algrenewcommand\algorithmicthen{\textbf{entonces}}
\algrenewcommand\algorithmicelse{\textbf{si no}}
\algrenewcommand\algorithmicfor{\textbf{para}}
\algrenewcommand\algorithmicforall{\textbf{para todo}}
\algrenewcommand\algorithmicdo{\textbf{hacer}}
\algrenewcommand\algorithmicwhile{\textbf{mientras}}
\algrenewcommand\algorithmicrepeat{\textbf{repetir}}
\algrenewcommand\algorithmicuntil{\textbf{hasta que}}
\algrenewcommand\algorithmicloop{\textbf{ciclo}}
\algrenewcommand\algorithmicprocedure{\textbf{procedimiento}}
\algrenewcommand\algorithmicfunction{\textbf{función}}
\algrenewcommand\algorithmicreturn{\textbf{devolver}}
\algrenewcomment[1]{\hfill{\color{comentario}\itshape // #1}}

% Las dos que algorithmicx no trae. El \ del final es obligatorio: sin
% el TeX se come el espacio y sale "hastan-1" en vez de "hasta n-1"
\algnewcommand\Hasta{\textbf{hasta}\ }
\algnewcommand\Intercambiar{\textbf{intercambiar}\ }

% --- abreviaturas que se usan a cada rato ---
\newcommand{\N}{\mathbb{N}}
\newcommand{\Z}{\mathbb{Z}}
\newcommand{\Q}{\mathbb{Q}}
\newcommand{\R}{\mathbb{R}}
\newcommand{\C}{\mathbb{C}}
\newcommand{\abs}[1]{\left\lvert #1 \right\rvert}
\newcommand{\norm}[1]{\left\lVert #1 \right\rVert}
\newcommand{\set}[1]{\left\{ #1 \right\}}
\newcommand{\Ord}[1]{\mathcal{O}\!\left(#1\right)}
% Mayuscula la notacion grande, minuscula la chica
\newcommand{\Omg}[1]{\Omega\!\left(#1\right)}
\newcommand{\Tht}[1]{\Theta\!\left(#1\right)}
\newcommand{\ord}[1]{o\!\left(#1\right)}
\newcommand{\omg}[1]{\omega\!\left(#1\right)}
% Los tres limites sobre n, que salen en cada renglon del criterio
% del cociente para la notacion asintotica
\newcommand{\LM}{\lim_{n \to \infty}}
\newcommand{\LS}{\limsup_{n \to \infty}}
\newcommand{\LI}{\liminf_{n \to \infty}}

% --- encabezado de cada sesión: título a la izquierda, fecha a la
%     derecha. El argumento corto es el que sale en el índice ---
\newcommand{\clase}[2]{%
  \section[#1]{#1 \hfill \normalfont\small\textit{#2}}%
}

% --- portadilla y tabla de contenidos. Se llama una vez, justo
%     despues de \begin{document}. El renglon del profesor va dentro
%     del center, asi que \profesor puede traer varios nombres
%     separados por \\ y salen apilados y centrados sin tocar nada ---
\newcommand{\portadilla}{%
  \begin{center}
    {\LARGE\bfseries\materia\par}
    \vspace{0.4em}
    {\large \profesor \par}
    \vspace{0.2em}
    {\small \periodo\quad-\quad\alumno\par}
    \vspace{0.3em}
    \rule{\linewidth}{0.5pt}
  \end{center}
  \vspace{0.5em}
  \tableofcontents
  \vspace{1em}
  \rule{\linewidth}{0.5pt}%
}

% =====================================================================
%  DOCUMENTOS DE SOLUCIONES
%  Lo que usa una tanda de ejercicios resueltos. Una libreta no toca
%  nada de aqui, pero vive en el preambulo compartido para no tener
%  dos archivos que arreglar cuando cambie el estilo.
% =====================================================================

% --- encabezado compacto: titulo, subtitulo y regla. Sustituye a
%     \portadilla cuando el documento es una tanda de ejercicios: sin
%     indice, porque se hojea de corrido, no se navega ---
\newcommand{\encabezadoSoluciones}[2]{%
  % El paquete algorithm se carga con la opcion [section] para que en
  % una libreta el pseudocodigo se numere por clase. Aqui no hay
  % secciones, asi que esa numeracion sale como "Algoritmo 0.1". Se
  % corrige al vuelo, y solo en los documentos que llaman a este
  % encabezado: la libreta no se entera.
  \renewcommand{\thealgorithm}{\arabic{algorithm}}%
  \begin{center}
    {\LARGE\bfseries #1\par}
    \vspace{0.3em}
    {\small\itshape #2\par}
    \vspace{0.5em}
    \rule{\linewidth}{0.5pt}
  \end{center}
  \vspace{0.3em}%
}

% --- hoja de un solo ejercicio: un renglon chico arriba con el
%     titulo y su numero, y ya. Sin titulo grande y sin pie de pagina,
%     porque en una hoja de dos paginas cada centimetro que se va en
%     adornos es un paso del desarrollo que no cabe ---
\newcommand{\encabezadoEjercicio}[3]{%
  % Misma correccion que en \encabezadoSoluciones: sin secciones, el
  % pseudocodigo saldria numerado "0.1"
  \renewcommand{\thealgorithm}{\arabic{algorithm}}%
  \fancyhf{}%
  \fancyhead[L]{\small\textsc{#1} (ejercicio #2 de #3)}%
  % El nombre va a la derecha porque las hojas de respuestas se
  % entregan y el profesor pide que cada una lo lleve
  \fancyhead[R]{\small\alumno}%
  \renewcommand{\headrulewidth}{0.4pt}%
  \pagestyle{fancy}%
  \thispagestyle{fancy}%
}

% --- bloque tematico, numerado con romanos: I, II, III. Agrupa los
%     ejercicios que comparten tema. No dibuja regla: la del
%     \problema que sigue ya separa el titulo del primer ejercicio,
%     y dos rayas seguidas dejan un hueco feo ---
\newcounter{parte}
\renewcommand{\theparte}{\Roman{parte}}
\newcommand{\parte}[1]{%
  \bigskip
  \refstepcounter{parte}%
  \noindent{\large\bfseries \theparte. #1\par}%
}

% --- un ejercicio resuelto. Se numera solo y corrido a lo largo del
%     documento, sin reiniciar en cada parte, para poder decir "el 14"
%     sin aclarar de que bloque ---
\newcounter{problema}
\newcommand{\problema}[1]{%
  \bigskip
  \noindent\rule{\linewidth}{0.4pt}\par
  \vspace{0.3em}
  \refstepcounter{problema}%
  \noindent{\bfseries Ejercicio \theproblema\ --- #1\par}
  \vspace{0.3em}%
}

% --- rotulo del paso que sigue: Caso base, Hipotesis inductiva, Paso
%     inductivo. El texto sigue en el mismo renglon ---
\newcommand{\paso}[1]{\par\vspace{0.3em}\noindent\textbf{#1}\ }

% --- la respuesta de un ejercicio, centrada y en un cuadro. El
%     argumento va en modo matematico, sin los delimitadores ---
\newcommand{\respuesta}[1]{%
  \[
    \boxed{\;#1\;}
  \]
}

% --- resaltado inline, como pasarle el plumon a media frase. No
%     admite matematicas adentro: soul rompe con $...$, asi que una
%     formula que hay que destacar va en \respuesta ---
\sethlcolor{marcador}
\newcommand{\resaltar}[1]{\hl{#1}}

% --- justificacion al final de un renglon de \align, en chiquito y
%     entre parentesis, para decir de donde salio el paso ---
\newcommand{\razon}[1]{\quad\text{\small\itshape (#1)}}

%############################ END OF PREAMBULO.TEX #############################
%###############################################################################

%  \begin{document}
%  \portadilla
%
%  Con dos o mas profesores se declara UN solo \profesor y los nombres
%  se separan con \\, uno por renglon:
%
%  \newcommand{\profesor}{Dra. Fulana de Tal \\
%                         Dr. Mengano de Cual}
% =====================================================================


% ==================== IDIOMA Y CODIFICACIÓN ====================
\usepackage[utf8]{inputenc}
\usepackage[T1]{fontenc}
\usepackage[spanish,es-tabla,es-noshorthands]{babel}
\usepackage{lmodern}
\usepackage{microtype}

% ==================== PÁGINA ====================
\usepackage[letterpaper,margin=2.3cm,headheight=14pt]{geometry}
\usepackage{parskip}
\usepackage{fancyhdr}
\usepackage{enumitem}

% ==================== MATEMÁTICAS ====================
\usepackage{amsmath,amssymb,amsthm,mathtools}

% ==================== FIGURAS, TABLAS, CÓDIGO ====================
\usepackage{graphicx}
\usepackage{booktabs}
\usepackage{array}
% float da el especificador [H], que clava una tabla donde se
% escribio. Una traza que LaTeX manda dos paginas adelante deja
% de leerse junto al paso que la produjo
\usepackage{float}
\usepackage{xcolor}
% tikz ya viene arrastrado por tcolorbox[most], pero se carga
% explicito para no depender de ese detalle; la libreria babel
% protege a tikz de los caracteres activos del español
\usepackage{tikz}
\usetikzlibrary{babel}
% pgfplots dibuja graficas de datos (ejes, escalas log) sobre tikz
\usepackage{pgfplots}
\pgfplotsset{compat=1.18}
\usepackage{listings}
\usepackage[ruled,section]{algorithm}
\usepackage{algpseudocode}
\usepackage[most]{tcolorbox}
% soul da \hl, el unico resaltado que parte renglon; \colorbox no
\usepackage{soul}
\usepackage{hyperref}

% =====================================================================
%  ESTILO
% =====================================================================

% --- encabezado y pie de cada página ---
\pagestyle{fancy}
\fancyhf{}
\fancyhead[L]{\small\textsc{\materia}}
\fancyhead[R]{\small\periodo}
\fancyfoot[C]{\small\thepage}
\renewcommand{\headrulewidth}{0.4pt}

% --- enlaces ---
\hypersetup{
  colorlinks=true,
  linkcolor=black,
  urlcolor=blue,
  citecolor=black,
  pdftitle={Apuntes de \materia},
  pdfauthor={\alumno}
}

% --- paleta ---
\definecolor{acento}{HTML}{1F4E79}
\definecolor{fondoIdea}{HTML}{EAF2FB}
\definecolor{fondoOjo}{HTML}{FDF0E6}
\definecolor{fondoPend}{HTML}{F2F0F7}
\definecolor{comentario}{HTML}{5A7D5A}
\definecolor{cadena}{HTML}{A03030}
\definecolor{fondoCodigo}{HTML}{F7F7F7}
\definecolor{fondoResultado}{HTML}{EAF5EC}
\definecolor{marcador}{HTML}{FFF3A3}

% --- entornos de teorema; numeran por sección, o sea por clase ---
\theoremstyle{definition}
\newtheorem{definicion}{Definición}[section]
\newtheorem{ejemplo}[definicion]{Ejemplo}
\newtheorem{ejercicio}[definicion]{Ejercicio}

\theoremstyle{plain}
\newtheorem{teorema}[definicion]{Teorema}
\newtheorem{lema}[definicion]{Lema}
\newtheorem{corolario}[definicion]{Corolario}
\newtheorem{proposicion}[definicion]{Proposición}

\theoremstyle{remark}
\newtheorem*{observacion}{Observación}
\newtheorem*{quick_sum}{Resumen rápido}

% --- cajas de color para lo que se relee antes del examen ---
\newtcolorbox{idea}[1][]{
  colback=fondoIdea, colframe=acento, boxrule=0.6pt,
  arc=2pt, left=6pt, right=6pt, top=5pt, bottom=5pt,
  title={Idea clave}, fonttitle=\bfseries\small, #1
}
\newtcolorbox{ojo}[1][]{
  colback=fondoOjo, colframe=orange!70!black, boxrule=0.6pt,
  arc=2pt, left=6pt, right=6pt, top=5pt, bottom=5pt,
  title={Ojito...}, fonttitle=\bfseries\small, #1
}
\newtcolorbox{pendiente}[1][]{
  colback=fondoPend, colframe=violet!60!black, boxrule=0.6pt,
  arc=2pt, left=6pt, right=6pt, top=5pt, bottom=5pt,
  title={Pendiente por resolver}, fonttitle=\bfseries\small, #1
}
\newtcolorbox{idear}[1][]{
  colback=fondoPend, colframe=pink!60!black, boxrule=0.6pt,
  arc=2pt, left=6pt, right=6pt, top=5pt, bottom=5pt,
  title={Idea Random}, fonttitle=\bfseries\small, #1
}

% --- el resultado al que llego un ejercicio, para que se vea de
%     lejos al hojear la hoja antes del examen ---
\newtcolorbox{resultado}[1][]{
  colback=fondoResultado, colframe=green!45!black, boxrule=0.6pt,
  arc=2pt, left=6pt, right=6pt, top=5pt, bottom=5pt,
  title={Resultado}, fonttitle=\bfseries\small, #1
}

% --- código ---
\lstdefinestyle{codigo}{
  backgroundcolor=\color{fondoCodigo},
  basicstyle=\ttfamily\footnotesize,
  keywordstyle=\color{acento}\bfseries,
  commentstyle=\color{comentario}\itshape,
  stringstyle=\color{cadena},
  numbers=left, numberstyle=\tiny\color{gray}, numbersep=8pt,
  frame=single, rulecolor=\color{gray!40},
  breaklines=true, showstringspaces=false,
  tabsize=4, captionpos=b,
  extendedchars=true, inputencoding=utf8,
  literate={á}{{\'a}}1 {é}{{\'e}}1 {í}{{\'i}}1
           {ó}{{\'o}}1 {ú}{{\'u}}1 {ñ}{{\~n}}1
           {Ü}{{\"U}}1 {ü}{{\"u}}1 {¿}{{?`}}1
           {¡}{{!`}}1
}
\lstset{style=codigo}
\renewcommand{\lstlistingname}{Código}
\renewcommand{\lstlistlistingname}{Índice de código}
% --- pseudocódigo; algorithmicx arma el cuerpo y algorithm el
%     flotante que lo numera y lo pone en el índice. Las palabras
%     clave vienen en inglés de fábrica y babel no las toca, así que
%     se renombran una por una ---
\floatname{algorithm}{Algoritmo}
\renewcommand{\listalgorithmname}{Índice de algoritmos}

\algrenewcommand\algorithmicrequire{\textbf{Entrada:}}
\algrenewcommand\algorithmicensure{\textbf{Salida:}}
\algrenewcommand\algorithmicend{\textbf{fin}}
\algrenewcommand\algorithmicif{\textbf{si}}
\algrenewcommand\algorithmicthen{\textbf{entonces}}
\algrenewcommand\algorithmicelse{\textbf{si no}}
\algrenewcommand\algorithmicfor{\textbf{para}}
\algrenewcommand\algorithmicforall{\textbf{para todo}}
\algrenewcommand\algorithmicdo{\textbf{hacer}}
\algrenewcommand\algorithmicwhile{\textbf{mientras}}
\algrenewcommand\algorithmicrepeat{\textbf{repetir}}
\algrenewcommand\algorithmicuntil{\textbf{hasta que}}
\algrenewcommand\algorithmicloop{\textbf{ciclo}}
\algrenewcommand\algorithmicprocedure{\textbf{procedimiento}}
\algrenewcommand\algorithmicfunction{\textbf{función}}
\algrenewcommand\algorithmicreturn{\textbf{devolver}}
\algrenewcomment[1]{\hfill{\color{comentario}\itshape // #1}}

% Las dos que algorithmicx no trae. El \ del final es obligatorio: sin
% el TeX se come el espacio y sale "hastan-1" en vez de "hasta n-1"
\algnewcommand\Hasta{\textbf{hasta}\ }
\algnewcommand\Intercambiar{\textbf{intercambiar}\ }

% --- abreviaturas que se usan a cada rato ---
\newcommand{\N}{\mathbb{N}}
\newcommand{\Z}{\mathbb{Z}}
\newcommand{\Q}{\mathbb{Q}}
\newcommand{\R}{\mathbb{R}}
\newcommand{\C}{\mathbb{C}}
\newcommand{\abs}[1]{\left\lvert #1 \right\rvert}
\newcommand{\norm}[1]{\left\lVert #1 \right\rVert}
\newcommand{\set}[1]{\left\{ #1 \right\}}
\newcommand{\Ord}[1]{\mathcal{O}\!\left(#1\right)}
% Mayuscula la notacion grande, minuscula la chica
\newcommand{\Omg}[1]{\Omega\!\left(#1\right)}
\newcommand{\Tht}[1]{\Theta\!\left(#1\right)}
\newcommand{\ord}[1]{o\!\left(#1\right)}
\newcommand{\omg}[1]{\omega\!\left(#1\right)}
% Los tres limites sobre n, que salen en cada renglon del criterio
% del cociente para la notacion asintotica
\newcommand{\LM}{\lim_{n \to \infty}}
\newcommand{\LS}{\limsup_{n \to \infty}}
\newcommand{\LI}{\liminf_{n \to \infty}}

% --- encabezado de cada sesión: título a la izquierda, fecha a la
%     derecha. El argumento corto es el que sale en el índice ---
\newcommand{\clase}[2]{%
  \section[#1]{#1 \hfill \normalfont\small\textit{#2}}%
}

% --- portadilla y tabla de contenidos. Se llama una vez, justo
%     despues de \begin{document}. El renglon del profesor va dentro
%     del center, asi que \profesor puede traer varios nombres
%     separados por \\ y salen apilados y centrados sin tocar nada ---
\newcommand{\portadilla}{%
  \begin{center}
    {\LARGE\bfseries\materia\par}
    \vspace{0.4em}
    {\large \profesor \par}
    \vspace{0.2em}
    {\small \periodo\quad-\quad\alumno\par}
    \vspace{0.3em}
    \rule{\linewidth}{0.5pt}
  \end{center}
  \vspace{0.5em}
  \tableofcontents
  \vspace{1em}
  \rule{\linewidth}{0.5pt}%
}

% =====================================================================
%  DOCUMENTOS DE SOLUCIONES
%  Lo que usa una tanda de ejercicios resueltos. Una libreta no toca
%  nada de aqui, pero vive en el preambulo compartido para no tener
%  dos archivos que arreglar cuando cambie el estilo.
% =====================================================================

% --- encabezado compacto: titulo, subtitulo y regla. Sustituye a
%     \portadilla cuando el documento es una tanda de ejercicios: sin
%     indice, porque se hojea de corrido, no se navega ---
\newcommand{\encabezadoSoluciones}[2]{%
  % El paquete algorithm se carga con la opcion [section] para que en
  % una libreta el pseudocodigo se numere por clase. Aqui no hay
  % secciones, asi que esa numeracion sale como "Algoritmo 0.1". Se
  % corrige al vuelo, y solo en los documentos que llaman a este
  % encabezado: la libreta no se entera.
  \renewcommand{\thealgorithm}{\arabic{algorithm}}%
  \begin{center}
    {\LARGE\bfseries #1\par}
    \vspace{0.3em}
    {\small\itshape #2\par}
    \vspace{0.5em}
    \rule{\linewidth}{0.5pt}
  \end{center}
  \vspace{0.3em}%
}

% --- hoja de un solo ejercicio: un renglon chico arriba con el
%     titulo y su numero, y ya. Sin titulo grande y sin pie de pagina,
%     porque en una hoja de dos paginas cada centimetro que se va en
%     adornos es un paso del desarrollo que no cabe ---
\newcommand{\encabezadoEjercicio}[3]{%
  % Misma correccion que en \encabezadoSoluciones: sin secciones, el
  % pseudocodigo saldria numerado "0.1"
  \renewcommand{\thealgorithm}{\arabic{algorithm}}%
  \fancyhf{}%
  \fancyhead[L]{\small\textsc{#1} (ejercicio #2 de #3)}%
  % El nombre va a la derecha porque las hojas de respuestas se
  % entregan y el profesor pide que cada una lo lleve
  \fancyhead[R]{\small\alumno}%
  \renewcommand{\headrulewidth}{0.4pt}%
  \pagestyle{fancy}%
  \thispagestyle{fancy}%
}

% --- bloque tematico, numerado con romanos: I, II, III. Agrupa los
%     ejercicios que comparten tema. No dibuja regla: la del
%     \problema que sigue ya separa el titulo del primer ejercicio,
%     y dos rayas seguidas dejan un hueco feo ---
\newcounter{parte}
\renewcommand{\theparte}{\Roman{parte}}
\newcommand{\parte}[1]{%
  \bigskip
  \refstepcounter{parte}%
  \noindent{\large\bfseries \theparte. #1\par}%
}

% --- un ejercicio resuelto. Se numera solo y corrido a lo largo del
%     documento, sin reiniciar en cada parte, para poder decir "el 14"
%     sin aclarar de que bloque ---
\newcounter{problema}
\newcommand{\problema}[1]{%
  \bigskip
  \noindent\rule{\linewidth}{0.4pt}\par
  \vspace{0.3em}
  \refstepcounter{problema}%
  \noindent{\bfseries Ejercicio \theproblema\ --- #1\par}
  \vspace{0.3em}%
}

% --- rotulo del paso que sigue: Caso base, Hipotesis inductiva, Paso
%     inductivo. El texto sigue en el mismo renglon ---
\newcommand{\paso}[1]{\par\vspace{0.3em}\noindent\textbf{#1}\ }

% --- la respuesta de un ejercicio, centrada y en un cuadro. El
%     argumento va en modo matematico, sin los delimitadores ---
\newcommand{\respuesta}[1]{%
  \[
    \boxed{\;#1\;}
  \]
}

% --- resaltado inline, como pasarle el plumon a media frase. No
%     admite matematicas adentro: soul rompe con $...$, asi que una
%     formula que hay que destacar va en \respuesta ---
\sethlcolor{marcador}
\newcommand{\resaltar}[1]{\hl{#1}}

% --- justificacion al final de un renglon de \align, en chiquito y
%     entre parentesis, para decir de donde salio el paso ---
\newcommand{\razon}[1]{\quad\text{\small\itshape (#1)}}

%############################ END OF PREAMBULO.TEX #############################
%###############################################################################


\begin{document}

\encabezadoEjercicio
  {Corrección de MAX-HEAPIFY por inducción sobre la altura}{20}{20}

\noindent\textbf{Convenciones.}
Se indexa desde $1$, con $\mathrm{padre}(i) = \lfloor i/2
\rfloor$, $\mathrm{izq}(i) = 2i$ y $\mathrm{der}(i) = 2i + 1$.

\noindent Demuestre que \textsc{Max-Heapify}$(A, i)$ restaura
correctamente la propiedad de max-heap en el subárbol con raíz $i$,
bajo la precondición de que los subárboles izquierdo y derecho de $i$
ya son max-heaps. Organice la demostración como una inducción sobre la
altura $h$ del subárbol: caso base $h = 0$; hipótesis inductiva, el
algoritmo funciona para subárboles de altura menor que $h$; paso
inductivo, analice qué sucede si la raíz ya es el mayor elemento y qué
sucede si debe intercambiarse con uno de sus hijos.

\paso{Notación y precondición:} sea $n$ el número de nodos,
$\ell = \mathrm{izq}(i)$ y $r = \mathrm{der}(i)$. El subárbol con raíz
$i$ es max-heap cuando $A[j] \ge A[c]$ para todo nodo $j$ de ese
subárbol y todo hijo $c$ de $j$ que exista. La precondición es que los
subárboles con raíz $\ell$ y $r$ ya lo cumplen, así que las únicas
parejas en duda son $(i, \ell)$ y $(i, r)$. El algoritmo toma el
índice $m$ del mayor entre $A[i]$, $A[\ell]$ y $A[r]$, contando solo
los hijos que existen; si $m = i$ se detiene, y si $m \ne i$
intercambia $A[i]$ con $A[m]$ y se llama sobre $m$.

Se prueba por inducción sobre $h$ la afirmación con sus dos cláusulas:
al terminar la llamada, (i) el subárbol con raíz $i$ es max-heap y
(ii) ninguna posición fuera de ese subárbol cambió. La cláusula (ii)
va dentro de la inducción porque el paso inductivo la necesita para
que la llamada recursiva no deshaga lo ya arreglado.

\paso{Caso base:} $h = 0$ significa $\ell = 2i > n$, o sea que $i$ no
tiene hijos; con $n = 10$ e $i = 7$, $\ell = 14 > 10$. El subárbol es
el solo nodo $i$ y no contiene ninguna pareja padre-hijo, de modo que
la condición $A[j] \ge A[c]$ se cumple por vacuidad. El algoritmo no
halla hijos con que comparar, deja $m = i$, no intercambia y termina,
así que (ii) también vale.

\paso{Hipótesis inductiva:} se supone la afirmación para todo nodo
cuyo subárbol tenga altura menor que $h$: si sus dos subárboles hijos
son max-heaps, la llamada lo deja como max-heap y no escribe en
ninguna posición fuera de ese subárbol.

\paso{Paso inductivo, caso 1, la raíz ya es el mayor:} aquí $m = i$, y
por como se eligió $m$,
\begin{align*}
  A[i] &\ge A[\ell]
         \razon{$A[i]$ ganó la comparación con el hijo izquierdo} \\
  A[i] &\ge A[r]
         \razon{y la del derecho, si $r$ existe}
\end{align*}
El algoritmo no escribe nada, luego (ii) vale y $A$ sigue siendo el de
la precondición. Las parejas padre-hijo del subárbol con raíz $i$ son
las dos de arriba y las que quedan dentro de los subárboles de $\ell$
y de $r$, que valen por la precondición. No queda ninguna otra pareja,
así que el subárbol con raíz $i$ es max-heap.

\paso{Paso inductivo, caso 2, hay que intercambiar:} aquí $m \ne i$.
Tómese $m = \ell$; con $m = r$ el argumento corre renglón por renglón
igual, porque abajo solo se usa que $m$ es hijo de $i$ y que su
subárbol era max-heap, nunca de qué lado cuelga. Sea $A'$ el arreglo
justo después del intercambio, o sea $A'[i] = A[m]$, $A'[m] = A[i]$ y
$A'[j] = A[j]$ en toda otra posición.

La nueva raíz domina al hijo que recibió el valor viejo:
\begin{align*}
  A'[i] &= A[m]
          \razon{el intercambio sube $A[m]$ a la raíz} \\
        &= \max\set{A[i],\ A[\ell],\ A[r]}
          \razon{$m$ es el índice del mayor de los tres} \\
        &\ge A[i]
          \razon{el máximo no es menor que ninguno} \\
        &= A'[m]
          \razon{el viejo $A[i]$ bajó a la posición $m$} \\
        &= A'[\ell]
          \razon{en este caso $m = \ell$}
\end{align*}
y también al otro hijo, que no se movió:
\begin{align*}
  A'[i] &= \max\set{A[i],\ A[\ell],\ A[r]}
          \razon{renglón anterior} \\
        &\ge A[r]
          \razon{el máximo no es menor que ninguno} \\
        &= A'[r]
          \razon{el intercambio solo escribió en $i$ y en $m$}
\end{align*}
Más aún, $A'[i]$ domina todo el subárbol de $m$ y no solo su raíz:
\begin{align*}
  A'[i] &= A[m]
          \razon{primer renglón de arriba} \\
        &\ge A[j] \quad \text{para todo $j$ del subárbol de $m$}
          \razon{encadenando $A[\text{padre}] \ge A[\text{hijo}]$ de
                 $m$ a $j$}
\end{align*}
y ese subárbol conserva en $A'$ los mismos valores salvo $A[m]$, que
fue sustituido por $A[i] \le A'[i]$. Todo valor que el subárbol de $m$
guarda en $A'$ es entonces a lo más $A'[i]$.

\paso{Dónde puede quedar la violación:} el intercambio escribió solo
en $i$ y en $m$, así que una pareja padre-hijo del subárbol de $i$
conserva sus dos valores salvo que toque una de esas dos posiciones.
Separando los cuatro grupos:
\begin{itemize}
  \item $(i, \ell)$ y $(i, r)$: recién probadas en $A'$.
  \item las parejas internas del subárbol de $r$: ninguno de sus nodos
        es $i$ ni $m$, conservan sus valores y valían por la
        precondición.
  \item las parejas internas del subárbol de $m$ que no tienen a $m$
        de padre: tampoco se tocaron, y valían porque el subárbol de
        $m$ era max-heap.
  \item $(m, \mathrm{izq}(m))$ y $(m, \mathrm{der}(m))$: las únicas
        que pueden fallar, porque $A'[m] = A[i]$ es un valor nuevo en
        esa posición mientras sus hijos siguen con el suyo.
\end{itemize}
Toda violación cae dentro del subárbol con raíz $m$, que es justo
donde el algoritmo vuelve a llamarse.

\paso{La llamada sobre $m$ cumple la precondición y baja de altura:}
los subárboles con raíz $\mathrm{izq}(m)$ y $\mathrm{der}(m)$ no
contienen a $i$, que es el padre de $m$, ni a $m$, que es su padre, de
modo que el intercambio no los tocó; y eran max-heaps por ser
subárboles del de $m$, que lo era. Por otro lado $m$ es hijo de $i$,
así que su subárbol es parte propia del de $i$ y su altura es a lo más
$h - 1 < h$. No tiene por qué ser exactamente $h - 1$, ya que el
subárbol derecho puede quedar un nivel más corto; por eso la hipótesis
se enuncia para toda altura menor que $h$.

\paso{Cierre:} por hipótesis inductiva la llamada deja el subárbol de
$m$ como max-heap y no escribe fuera de él. Las parejas internas de
ese subárbol quedan bien por (i) de la hipótesis; el subárbol de $r$ y
la pareja $(i, r)$ no volvieron a tocarse por (ii); y en la pareja
$(i, m)$ el valor de $i$ no cambió, mientras que la llamada solo
permuta valores del subárbol de $m$, todos a lo más $A'[i]$, así que
el valor final en $m$ tampoco supera al de $i$. Todas las parejas del
subárbol de $i$ se cumplen, que es (i). Lo único escrito fue el
intercambio en $i$ y $m$ más lo que hizo la llamada dentro del
subárbol de $m$, todo dentro del subárbol de $i$, que es (ii). La
inducción queda cerrada.
\respuesta{A[j] \ge A[c] \quad \text{para todo $j$ del subárbol con
  raíz $i$ y todo hijo $c$ de $j$}}

\end{document}

%################# END OF EJERCICIO20CORRECCIONMAXHEAPIFY.TEX ##################
%###############################################################################
